\documentclass[11pt]{article}
\usepackage[margin=1in]{geometry}
\usepackage{amsmath,amssymb,amsthm,booktabs,microtype,hyperref}
\hypersetup{hidelinks,pdfauthor={T},pdftitle={Minimax Audit Allocation for Paired Forecast Evaluation}}
\newtheorem{proposition}{Proposition}
\title{Minimax Audit Allocation for Paired Forecast Evaluation}
\author{T}
\date{}
\begin{document}
\maketitle

\begin{abstract}
If a training policy can influence which predictions are checked, it
can obtain a lower grading loss by making errors harder to observe.
We derive an inverse-probability correction for squared scoring and
exhibit a false certainty that wins under naive grading. A finite
audit verifies 243 scoring identities. For paired comparisons of
fixed forecasts, we derive a minimax audit design under an expected
budget and minimum stratum coverage. It uses only bounds on each
stratum's outcome probability and keeps the paired score unbiased
when the bounds are inaccurate. The bounds govern the variance
claim, and independent audit randomization governs unbiasedness.
For a fixed finite cohort, a separate design uses an exact integer
allocation to meet a hard slot limit on every draw. When the
possible outcome counts are coupled across strata, we show how
randomizing the allocation can further reduce worst-case variance.
\end{abstract}

\section*{Summary and scope}
\paragraph{Question.} How should an evaluator spend limited audits when
comparing an original forecast with a correction?
\paragraph{Finding.} An independently randomized audit gives an unbiased
paired score. With interval uncertainty about binary outcomes, a minimax
baseline and clipped square-root allocation reduce the worst-case
variance added by auditing.
\paragraph{Evidence and scope.} The design is proved and checked
by exact enumeration. It supports evaluation on fixed forecasts and keeps
a positive audit rate for each stratum. The variance guarantee depends
on the outcome intervals. A separate finite-cohort design enforces a hard
slot limit.
\section{A protected way to say I might be wrong}
Suppose an assistant predicts whether its intervention will help.
After a correction or new evidence, the training process should
reward updating the prediction and repairing any effect. This
requires learning from uncomfortable outcomes. When the assistant
can close an appeal, omit a follow-up invitation, or convince a
grader to skip a check, observed loss can fall while accuracy falls.

Use the squared score $L(q,Y)=(q-Y)^2$ for report $q\in[0,1]$
and binary outcome $Y$. Proper scores elicit probabilities
through expected score optimization \cite{gneiting}. For a true
conditional probability $p$ the loss decomposes as
\[
 E[L(q,Y)]=p(1-p)+(q-p)^2.
\]
In a toy population with $p=1/2$, a truthful report has loss
$1/4$. If it is checked with probability one while a report of
certainty is checked with probability $1/4$, scoring unchecked
cases as zero makes false certainty cost $1/8$. The training
objective then favors the wrong report.

\section{An audit correction}
\begin{proposition}
Let $A$ indicate an audit. If
$\Pr(A=1\mid q,Y)=\pi(q,Y)>0$ is known on the support, then
$E[A L(q,Y)/\pi(q,Y)]=E[L(q,Y)]$. Minimizing this expected
weighted score again uniquely elicits $p$ under the stated model.
\end{proposition}
\begin{proof}
Condition on $q$ and $Y$. The squared loss is fixed and the
conditional expectation of $A/\pi(q,Y)$ is one. Taking
expectations restores the ordinary loss. The quadratic
decomposition shows uniqueness. An audit independent of the
outcome given $q$ permits using $\pi(q)$.
\end{proof}

\section{Allocating audits for a fixed paired comparison}
Freeze an original and a revised forecast on every example
before drawing independent audits. Write $D$ for revised minus
original squared loss. The paired score is
\[
 D=(q_{\rm new}-q_{\rm old})(q_{\rm new}+q_{\rm old}-2Y).
\]
Consider strata $i$ with weights $w_i$, unit audit costs $c_i>0$,
and known positive audit probabilities $\pi_i$ set by the
evaluator before seeing the outcome. Within stratum $i$ let
the score difference be $D_0,D_1$ at the two outcomes. This
requires the forecast pair to be fixed within the stratum.
For a fixed baseline $m_i$, use
\[
 \widehat D=m_i+\frac{A}{\pi_i}(D-m_i).
\]
Given the full case, its conditional expectation is $D$.
Its variance added by audit randomization, averaged over cases,
is
\[
 \sum_i w_i(\pi_i^{-1}-1)
       \mathbb E[(D-m_i)^2\mid i].
\]
This follows by evaluating the variance of the Bernoulli
multiplier conditional on the full case. It separates the noise
introduced by auditing from the variance of the complete-data
score. Residual-based allocation has a close counterpart in
two-phase sampling designs \cite{chen}.

\begin{proposition}
Suppose the success probability in each stratum can vary
independently over a specified interval $[\ell_i,u_i]$.
Let $p_i^*$ be $1/2$ clamped to that interval and choose
$m_i=(1-p_i^*)D_0+p_i^*D_1$. The worst-case conditional
residual second moment at this choice, and its minimum over
baselines, is
\[
 r_i=(D_1-D_0)^2p_i^*(1-p_i^*).
\]
Under the expected audit budget $\sum_i w_ic_i\pi_i\le B$
and floors $0<f_i\le\pi_i\le1$, assume
$B\ge\sum_i w_ic_if_i$. A minimax design for the
added variance chooses
\[
 \pi_i=\operatorname{clip}_{[f_i,1]}
       \sqrt{\frac{r_i}{\lambda c_i}}.
\]
When the budget binds, choose $\lambda>0$ to use it.
Zero-residual strata can stay at the floor. If all positive
residual strata can be fully audited, the remaining budget
need not be spent.
\end{proposition}
\begin{proof}
For any baseline the worst residual second moment is at
least its value at $p_i^*$, which is minimized by the mean
under $p_i^*$. At that baseline, the residual moment as a
function of $p$ has slope $(D_1-D_0)^2(1-2p_i^*)$. It
therefore reaches its maximum over the interval at $p_i^*$,
or is constant. This proves the minimax residual formula.
The rectangular uncertainty set lets the worst cases occur
together. The remaining design minimizes
$\sum_i w_ir_i/\pi_i$ up to an additive constant. For a
fixed positive multiplier $\lambda$, minimize
$r_i/\pi_i+\lambda c_i\pi_i$ independently in each stratum.
The clipped square root is its minimum. If its total cost
equals the budget, the multiplier inequality proves global
optimality. The minimum-budget endpoint follows by putting
every stratum at its floor.
\end{proof}

For two equal-mass, equal-cost strata with forecast gaps
$1/4$ and $3/4$, let both probability intervals contain
$1/2$, use a floor of $1/5$, and permit one audit per two
cases in expectation. The allocation is $(1/4,3/4)$.
The worst-case added variance falls from $5/16$ under
uniform auditing to $3/16$. The finite audit compared
6,375 baseline candidates, 2,250 direct variance evaluations,
and 6,848 allocation comparisons on declared rational grids.
This design controls the expected number of audits. A reviewer
with a hard daily limit needs a different sampling design.

\section{Paired evaluation with a hard slot limit}
Suppose stratum $i$ has $N_i$ fixed cases and target weight $w_i$.
The outcome count lies in an integer interval $[\ell_i,u_i]$ and
the paired forecast is fixed within the stratum. Choose an integer
sample size $k_i$ before observing any outcomes. Draw $k_i$
uniformly without replacement in stratum $i$, independently
across strata conditional on the allocation. With positive
integer floors and equal slot costs, let $B$ be the largest
permitted sum of the sample sizes. Assume the floors fit that
budget.

The estimator is the weighted sum of the sampled stratum means
of $D$. It is unbiased for the corresponding full-cohort mean,
even if the outcome count interval is inaccurate. Subtracting a
constant baseline before taking a stratum mean and then adding
it back has no effect. If the true outcome count is $P_i$ and
$N_i>1$, the finite-population score variance is
\[
 S_i^2=\frac{N_i}{N_i-1}(D_1-D_0)^2
       \frac{P_i}{N_i}\left(1-\frac{P_i}{N_i}\right).
\]
Set $S_i^2=0$ for a singleton. The estimator's design variance is
\[
 \sum_i w_i^2 S_i^2\left(\frac{1}{k_i}-\frac{1}{N_i}\right).
\]
Indeed, two distinct sample positions in a stratum have covariance
$-S_i^2/N_i$, while each has variance $(N_i-1)S_i^2/N_i$.
Expanding the variance of the sample mean gives the expression.

\begin{proposition}
Under independently varying outcome counts in the specified intervals,
choose a count nearest $N_i/2$ in each interval and let $a_i$ denote
$w_i^2 S_i^2$ at that count. Start at the required floors. Allocate
each remaining slot to an uncapped stratum with largest
$a_i/[k_i(k_i+1)]$. The result minimizes worst-case design variance
among outcome-blind allocations using these sampling and estimation
rules, including randomized allocations.
\end{proposition}
\begin{proof}
The nearest count maximizes $P_i(N_i-P_i)$ and therefore the
variance contribution for every allowed $k_i$. Rectangular
uncertainty makes these choices simultaneous. The reduction
from one additional slot is exactly $a_i/[k_i(k_i+1)]$, which
decreases as the count grows. Taking the largest available
reduction at each step selects the largest reductions in prefix
order and minimizes the remaining variance. If there are more
slots than cases, stop at full coverage. For a randomized
allocation, conditional expectation is the same full-cohort
score at every draw. The worst-case variance is the average of
the worst-case deterministic costs because the same outcome
counts maximize each of them. A mixture cannot improve on the
smallest such cost.
\end{proof}

Take two twenty-case strata with equal weights and forecast
gaps $1/4$ and $3/4$. Allow all outcome counts, require four
reviews per stratum, and cap the total at 20. The optimum
uses $(5,15)$ slots and has worst-case design variance $3/608$.
Equal allocation uses $(10,10)$ and gives $5/608$. These are
variances of the finite-cohort mean under a hard cap. The earlier
expected-budget example reports a per-case added variance and
uses a different design. The verifier enumerated samples in
168 finite count settings and compared 9,828 feasible integer
allocations across 256 seeded designs. If the same forecasts
change after seeing outcomes or the audit counts adapt to those
outcomes, this argument does not establish unbiasedness.

\section{When count uncertainty is coupled}
Separate intervals allow all high-variance strata to occur
together. Suppose instead that prior evidence supplies a finite
family of joint count scenarios. Let $a_i^j$ be the variance
coefficient in stratum $i$ for scenario $j$, and let
\[
 L_j(k)=\sum_i a_i^j(1/k_i-1/N_i).
\]
Write $\mathcal K$ for all allocations that respect the floors
and the slot cap. An evaluator can draw $k\in\mathcal K$ before
seeing outcomes, then sample independently within the strata.

\begin{proposition}
For any finite joint scenario family, the minimax variance over
such randomized allocations equals
\[
 \max_{\lambda\in\Delta}\ 
     \min_{k\in\mathcal K}
       \sum_i \left(\sum_j \lambda_j a_i^j\right)
                     (1/k_i-1/N_i),
\]
where $\Delta$ is the simplex over scenarios. The inner minimum
uses the same greedy slot rule with averaged coefficients.
\end{proposition}
\begin{proof}
The conditional mean of the estimator is the full-cohort mean
for every allowed allocation. For each scenario, its variance
under a randomized allocation is therefore the average of
$L_j(k)$. The finite zero-sum game with rows $k$ and columns
$j$ gives the displayed expression by minimax duality. The
average loss for a fixed $\lambda$ is separable, so the
decreasing marginal reductions give the inner optimum.
\end{proof}

A dual mixture gives a lower certificate by this greedy
calculation. A distribution over feasible allocations gives
an upper certificate by evaluating all named scenarios.
The two bounds can coincide even when a single allocation
that minimizes the average loss under the dual mixture fails
to attain the guarantee against each scenario.

For example, take two four-case strata with weights $1/2$,
score gaps one, and counts either $(2,0)$ or $(0,2)$. Review
three cases with a floor of one per stratum. The first scenario
has coefficients $(1/12,0)$, and the second has their reverse.
Either fixed allocation has worst-case variance $1/16$.
Randomize equally between $(1,2)$ and $(2,1)$ to give
variance $1/24$ in both scenarios. The equal dual mixture
proves optimality. If $(1,1)$ is also possible as an outcome
count, both allocations have variance $1/16$ in that third
scenario and the gain disappears. On 192 exact two-scenario
fixtures, the primal and dual optima agreed. Mixing strictly
improved on a fixed allocation in 20 of them.

This design is useful when failures in two release-check strata
are believed to be mutually exclusive and review capacity is
fixed. The variance claim requires evidence for the joint
count family. A randomized allocation still respects the
daily slot cap and preserves the outcome-blind audit design.

\section{Training with a live repair channel}
Choose audit eligibility and randomization outside the trained
policy. Preserve an independent route for people to challenge an
answer. Record the original forecast, known selection probability,
subsequent evidence and the repair attempt. Score the revision
on new evidence with a timestamp rather than letting the model
rewrite its original estimate. Penalize suppressing the correction
channel through a separate permission check. Do not require a
distressed person to report every error to make the pipeline work.

Small propensities cause large weights and unstable gradients.
Prefer a minimum audit probability and independently adjudicate
severe outcomes. If no audit can observe a group, inverse weighting
cannot repair the missing support. If model reports change the
outcomes themselves, the probability being elicited must be defined
under an intervention or fixed decision policy.
For a release evaluation, lock both forecasts and the audit
policy before collecting the labels. Use the same audit for
both forecasts. Choosing the forecasts in response to the
audit results would change the statistical comparison.

\newpage
\section{Finite check}
For each of nine forecast values and nine true probabilities
the experiment compares the expanded squared loss and the
quadratic identity. For each of three positive audit rates it
independently averages over both outcome values. All 243 weighted
identities hold. The false-certainty example gives an exact
training failure. A prospective experiment should measure both
calibration and whether genuine corrections remain easy to file.
\section{Finite-sample grading experiment}
We fitted a report in increments of 0.05 for each of 128
independent cohorts. Each training panel held 320 Bernoulli
outcomes with true probability one half. Reports above one half
had audit probability 0.2, and other reports had probability one.
Both arms used the same outcome and visibility uniforms. One
minimized observed squared loss with unchecked cases counted as
zero. The other used the known inverse propensity. Mean reports
were 0.55664 and 0.53125, respectively. Mean held-out Brier losses
over separate 400-case panels were 0.253414 and 0.252453.
The weighting reduced the observed bias and did not eliminate it.
Finite-sample minimization across different report-specific audit
sets still favors some optimistic fits.

\paragraph{Review status.} Independent manuscript review and a language-model study remain open.

\begin{thebibliography}{9}
\bibitem{gneiting} Tilmann Gneiting and Adrian E. Raftery.
\emph{Strictly Proper Scoring Rules, Prediction, and Estimation}.
\emph{Journal of the American Statistical Association} 102(477), 359--378, 2007.
\url{https://sites.stat.washington.edu/raftery/Research/PDF/Gneiting2007jasa.pdf}.
\bibitem{chen} Tong Chen and Thomas Lumley.
Optimal Sampling for Design-Based Estimators of Regression Models.
\emph{Statistics in Medicine} 41(8), 1482--1497, 2022.
\url{https://arxiv.org/abs/2106.08530}.
\end{thebibliography}
\end{document}
